\documentclass{situationist-preprint}

\title{The Zombie's Own Testimony}
\msid{MND-05-b}
\msfamily{MND-05}
\msdiscipline{Philosophy of Mind}
\mstype{Companion manuscript}
\mskeywords{zombies; conceivability; phenomenal judgment; conceivability-possibility principle}

\begin{document}
\maketitle

\begin{abstract}
A zombie, being our physical duplicate, says what we say about consciousness: that it is conscious, that zombies are conceivable, and that physicalism is therefore false. We argue that the zombie's testimony gives an independent argument for the conclusion of MND-05. The zombie's conceivability argument has premises it finds as compelling as we find ours, and a conclusion that, in its world, is false. A conceivability-possibility principle that holds for propositions formulated with phenomenal concepts would make the zombie's argument sound. So the principle fails for such propositions, whether or not zombies are possible.
\end{abstract}

\section{The argument}

Consider a world containing a zombie twin of a philosopher who accepts the conceivability argument. Being physically identical to her, the zombie produces the same words and is in the same functional states: it asserts that it has experiences, that it can conceive of a physical duplicate lacking them, and that what is ideally conceivable is possible. It infers that physicalism is false of its world. \citet{Chalmers1996} grants that a zombie would make such judgments, and calls the fact that our judgments about consciousness are physically explicable the paradox of phenomenal judgment.

In the zombie's world, by hypothesis, physicalism about everything is true: the world contains only physical properties, and the zombie has no experiences. The zombie's conclusion is therefore false. Its argument has the same form as ours, and the zombie's grounds for each premise are, physically and functionally, our grounds. \citet{Balog1999} draws the moral that the conceivability argument proves too much.

\begin{proposition}\label{p:fail}
If zombies are conceivable, the conceivability-possibility principle fails for some propositions formulated with phenomenal or quasi-phenomenal concepts.
\end{proposition}

\begin{proof}
If zombies are conceivable, a zombie world is conceivable, and in it the zombie's argument has an ideally conceivable premise, a conceivability-possibility premise in the same form as ours, and a false conclusion. The conceivability premise is formulated with the concepts the zombie uses in place of phenomenal concepts, which are functionally identical to ours. If the principle held for propositions formulated with concepts of that functional kind, the zombie's argument would be sound and its conclusion true; it is false. So the principle fails for such propositions, at least in the zombie's world, and since it is offered as an a priori principle, it fails as such.
\end{proof}

The proposition is reached without assuming physicalism about the actual world, which MND-05-a assumes (Assumption~3.1), and without the account of phenomenal concepts as directly referring (Definition~2.1). It uses only the conceivability of zombies, which the argument's proponent grants. Its conclusion is what Theorem~3.3 and Corollary~3.4 of MND-05-a establish by another route: the conceivability-possibility principle fails, and the failure is confined to propositions formulated with concepts of the phenomenal kind. Combined with Lemma~3.2 of MND-05-a, which uses physicalism, it yields that zombies are conceivable but not possible.

\begin{objection}
The zombie has no phenomenal concepts, and its argument is not the same argument. Ours is sound because our concepts refer to experiences, and we are acquainted with them.
\end{objection}

\begin{reply}
The reply concedes that nothing available to the reasoner distinguishes the sound argument from the unsound one, since the zombie's grounds are functionally ours. An a priori principle whose correct application depends on a difference the reasoner cannot detect a priori is not a principle that conceivability supports. The proponent may hold that acquaintance makes the difference, but that is to rest the argument on acquaintance rather than on conceivability, which MND-05 does not dispute.
\end{reply}

\section{Conclusion}

The zombie's testimony matches ours and its conclusion is false. The conceivability-possibility principle, applied to propositions formulated with phenomenal concepts, cannot distinguish the two, and so fails for them. Zombies are conceivable, and their conceivability does not establish their possibility.

\begin{thebibliography}{9}\small
\bibitem[Balog(1999)]{Balog1999} Balog, K. (1999). Conceivability, possibility, and the mind-body problem. \emph{The Philosophical Review}, 108(4), 497--528.
\bibitem[Chalmers(1996)]{Chalmers1996} Chalmers, D.~J. (1996). \emph{The Conscious Mind: In Search of a Fundamental Theory}. Oxford University Press.
\end{thebibliography}

\end{document}
